\makeatletter
\def\input@path{{./}{../../../preamble/}{book2026-2027/preamble/}}
\makeatother
\input{preamble.tex}
\usepackage{needspace}
\hypersetup{pdftitle={Grade 10 Term 2 Discrete Probability}}
\setlength{\emergencystretch}{2em}
\title{Grade 10, Term 2, Lesson 2, Criterion C5\\Geometric Waiting for the First Success}
\author{}
\date{}
\begin{document}
\maketitle
\setcounter{tocdepth}{1}
\setcounter{secnumdepth}{2}
\phantomsection\label{toc}
\tableofcontents
\section*{Evaluated criterion}
\textbf{C5:} Models waiting-time situations with geometric distributions, calculating the probability of the first success occurring on a given trial in financial or business contexts.
\section{Introduction}\label{sec:t2-c5-introduction}
All data in this activity are hypothetical teaching models. Probabilities have no monetary units. Keep exact expressions where useful and round final probabilities to four decimal places.
\subsection{Wait for the first success}
A school enterprise contacts visitors one by one until somebody buys. Define success as a purchase and $T$ as the call number of the first purchase. The first call is trial 1, so $T$ has support $\{1,2,3,\ldots\}$ calls, provided $0<p<1$. Trials are assumed independent with constant purchase probability $p>0$ and may continue as needed.
\textbf{Check.} Can $T=0$? \emph{Suggested answer:} No. At least one trial is needed to observe a first success.


\subsection{The sequence before stopping}
For $T=4$, the sequence must be FFFS: three failures followed by success. Its probability is $(1-p)^3p$. There is only one such pattern because the first success fixes the final position. $T$ counts all trials; $F=T-1$ counts only failures. These are different variables with different supports.
\textbf{Check.} If $T=4$, how many failures occur? \emph{Suggested answer:} $F=T-1=3$.


\subsection{By and after}
By trial $m$ means $T\leq m$ and includes success on trial $m$. After trial $m$ means $T>m$, so the first $m$ trials all fail. A trial count is not elapsed time unless each trial duration is specified. A 2 USD call cost would give a separate total cost $C=2T$ USD.
\textbf{Check.} Does after trial 3 include $T=3$? \emph{Suggested answer:} No. It means $T=4,5,\ldots$.

\textbf{Check.} How is this different from binomial? \emph{Suggested answer:} Binomial fixes the number of trials and counts successes; geometric fixes one success and counts trials until it.

\section{General worked problems}
\Needspace{14\baselineskip}
\subsection{Introductory general problem: first success on trial three}
Independent trials have success probability $1/2$.

\textbf{Questions/Tasks}
\begin{enumerate}
\item Calculate $P(T=3)$. 
\end{enumerate}
\subsubsection*{Complete worked solution}
\begin{center}
\begin{tabular}{llp{0.48\linewidth}}
\toprule
Symbol & Value & Meaning in this problem \\ \midrule
$p$ & $0.5$ & Success probability per trial \\
$t$ & $3$ trials & Requested first-success trial \\
\bottomrule
\end{tabular}
\end{center}

\par\noindent\begin{minipage}{\linewidth}
\textbf{Step 1: Define the variable.} $T$ is the trial number of the first stated success, in trials, with support $\{1,2,\ldots\}$. $p=0.5$ and $1-p=0.5$.
\end{minipage}\par\medskip
\par\noindent\begin{minipage}{\linewidth}
\textbf{Step 2: Check assumptions.} The trials are binary, independent and have constant success probability; the process stops at the first success. Thus a geometric total-trial model applies.
\end{minipage}\par\medskip
\par\noindent\begin{minipage}{\linewidth}
\textbf{Step 3: Translate and calculate.} Use
\[P(T=t)=(1-p)^{t-1}p.\]
Then
\[P(T=3)=(0.5)^{2}(0.5)\approx 0.1250.\]
\end{minipage}\par\medskip
\par\noindent\begin{minipage}{\linewidth}
\textbf{Step 4: Check and interpret.} The result is in $[0,1]$; by-trial and after-trial probabilities sum to 1 because their events are complements. For the exact event, 2 failures precede success on trial 3. 
\end{minipage}\par\medskip
\subsection{Bridge to the general model}
\[
P(T=t)=(1-p)^{t-1}p,\quad t=1,2,\ldots;\qquad
P(T>m)=(1-p)^m,\quad P(T\leq m)=1-(1-p)^m.
\]
There are $t-1$ failures before the final success. Powers mean repeated multiplication. The infinite PMF sums to 1 for $0<p\leq1$; when $p=1$, $T=1$ certainly. If $p=0$, success never occurs and this is not a proper finite waiting-time distribution. In failure-count notation $F=T-1$, $P(F=f)=(1-p)^fp$, with support $\{0,1,\ldots\}$ for $0<p<1$.
\Needspace{14\baselineskip}
\subsection{Advanced general problem: a window and a changed starting point}
Trials are independent with success probability 0.3.

\textbf{Questions/Tasks}
\begin{enumerate}
\item Find $P(2<T\leq5)$. Given the first two trials failed, find the chance of a success in the next three.
\end{enumerate}
\subsubsection*{Complete worked solution}
\par\noindent\begin{minipage}{\linewidth}
\textbf{Step 1: Define and model.} $T$ is the total trial count until first success, support $\{1,2,\ldots\}$ trials. Constant $p=0.3$ and independence justify geometric waiting.
\end{minipage}\par\medskip
\par\noindent\begin{minipage}{\linewidth}
\textbf{Step 2: Translate.} $2<T\leq5$ includes trials 3, 4, 5. Its unconditional probability is $P(T>2)-P(T>5)$.
\end{minipage}\par\medskip
\par\noindent\begin{minipage}{\linewidth}
\textbf{Step 3: Calculate.} $P(2<T\leq5)=(1-p)^2-(1-p)^5=0.7^2-0.7^5=0.3219$. The conditional probability, written $P(T\leq5\mid T>2)$, means given no success in two trials. Independence leaves the next three trials unchanged: $1-(1-p)^3=1-0.7^3=0.6570$.
\end{minipage}\par\medskip
\par\noindent\begin{minipage}{\linewidth}
\textbf{Step 4: Check and interpret.} The unconditional window probability is $0.32193$; dividing by $P(T>2)=0.49$ gives $0.657$. Waiting already observed changes the information, so these two probabilities answer different questions.
\end{minipage}\par\medskip
\section{Notation reference and practice procedure}
\begin{center}
\begin{tabular}{lp{0.72\linewidth}}
\toprule
Symbol & Meaning and units \\ \midrule
$T,t$ & Total trials until first success and one possible value \\
$p$ & Constant success probability per trial \\
$m$ & Trial cutoff; a positive integer \\
$F=T-1$ & Failures before first success \\
$P(A\mid B)$ & Probability of A given that B is known \\
\bottomrule
\end{tabular}
\end{center}
\begin{enumerate}
\item Define the random variable in words and give its units and support.
\item List parameters and check that the model assumptions match the situation.
\item Translate ordinary language into an event, paying attention to endpoints.
\item Write the symbolic formula first, then substitute without rounding intermediate values.
\item Check bounds, total probability or a complementary calculation; interpret the answer in context.
\end{enumerate}
\Needspace{19\baselineskip}
\section{Foundational contextual practice}\label{sec:t2-c5-foundational}
\Needspace{14\baselineskip}
\subsection{First sale on call four}
Independent sales calls have constant purchase probability 0.2; success means a purchase.

\textbf{Questions/Tasks}
\begin{enumerate}
\item Calculate $P(T=4)$. 
\end{enumerate}
\subsubsection*{Complete worked solution}
\par\noindent\begin{minipage}{\linewidth}
\textbf{Step 1: Define the variable.} $T$ is the trial number of the first stated success, in trials, with support $\{1,2,\ldots\}$. $p=0.2$ and $1-p=0.8$.
\end{minipage}\par\medskip
\par\noindent\begin{minipage}{\linewidth}
\textbf{Step 2: Check assumptions.} The trials are binary, independent and have constant success probability; the process stops at the first success. Thus a geometric total-trial model applies.
\end{minipage}\par\medskip
\par\noindent\begin{minipage}{\linewidth}
\textbf{Step 3: Translate and calculate.} Use
\[P(T=t)=(1-p)^{t-1}p.\]
Then
\[P(T=4)=(0.8)^{3}(0.2)\approx 0.1024.\]
\end{minipage}\par\medskip
\par\noindent\begin{minipage}{\linewidth}
\textbf{Step 4: Check and interpret.} The result is in $[0,1]$; by-trial and after-trial probabilities sum to 1 because their events are complements. For the exact event, 3 failures precede success on trial 4. 
\end{minipage}\par\medskip
\secInicio[][sec:t2-c5-foundational]
\Needspace{14\baselineskip}
\subsection{Loan approval on the second application}
In a simplified hypothetical model, independent applications each have approval probability 0.4; success means approval. Stop at the first approval.

\textbf{Questions/Tasks}
\begin{enumerate}
\item Calculate $P(T=2)$. State the number of rejected applications on this event.
\end{enumerate}
\subsubsection*{Complete worked solution}
\par\noindent\begin{minipage}{\linewidth}
\textbf{Step 1: Define the variable.} $T$ is the trial number of the first stated success, in trials, with support $\{1,2,\ldots\}$. $p=0.4$ and $1-p=0.6$.
\end{minipage}\par\medskip
\par\noindent\begin{minipage}{\linewidth}
\textbf{Step 2: Check assumptions.} The trials are binary, independent and have constant success probability; the process stops at the first success. Thus a geometric total-trial model applies.
\end{minipage}\par\medskip
\par\noindent\begin{minipage}{\linewidth}
\textbf{Step 3: Translate and calculate.} Use
\[P(T=t)=(1-p)^{t-1}p.\]
Then
\[P(T=2)=(0.6)^{1}(0.4)\approx 0.2400.\]
\end{minipage}\par\medskip
\par\noindent\begin{minipage}{\linewidth}
\textbf{Step 4: Check and interpret.} The result is in $[0,1]$; by-trial and after-trial probabilities sum to 1 because their events are complements. For the exact event, 1 failures precede success on trial 2. On this event $F=T-1=1$ rejection. Actual repeat applications may not be independent; the result relies on the supplied hypothetical assumption.
\end{minipage}\par\medskip
\secInicio[][sec:t2-c5-foundational]
\Needspace{19\baselineskip}
\section{Intermediate contextual practice}\label{sec:t2-c5-intermediate}
\Needspace{14\baselineskip}
\subsection{Payment completed by attempt three}
Independent payment attempts each succeed with probability 0.7. Stop when payment completes.

\textbf{Questions/Tasks}
\begin{enumerate}
\item Calculate $P(T\leq3)$. 
\end{enumerate}
\subsubsection*{Complete worked solution}
\par\noindent\begin{minipage}{\linewidth}
\textbf{Step 1: Define the variable.} $T$ is the trial number of the first stated success, in trials, with support $\{1,2,\ldots\}$. $p=0.7$ and $1-p=0.3$.
\end{minipage}\par\medskip
\par\noindent\begin{minipage}{\linewidth}
\textbf{Step 2: Check assumptions.} The trials are binary, independent and have constant success probability; the process stops at the first success. Thus a geometric total-trial model applies.
\end{minipage}\par\medskip
\par\noindent\begin{minipage}{\linewidth}
\textbf{Step 3: Translate and calculate.} Use
\[P(T\leq m)=1-(1-p)^m.\]
Then
\[P(T\leq3)=1-(0.3)^3\approx 0.9730.\]
\end{minipage}\par\medskip
\par\noindent\begin{minipage}{\linewidth}
\textbf{Step 4: Check and interpret.} The result is in $[0,1]$; by-trial and after-trial probabilities sum to 1 because their events are complements. The endpoint is included for by and excluded for after. 
\end{minipage}\par\medskip
\secInicio[][sec:t2-c5-intermediate]
\Needspace{14\baselineskip}
\subsection{Still waiting after five enquiries}
Independent enquiries each lead to a signed contract with probability 0.15. A contract is success.

\textbf{Questions/Tasks}
\begin{enumerate}
\item Calculate $P(T>5)$. 
\end{enumerate}
\subsubsection*{Complete worked solution}
\par\noindent\begin{minipage}{\linewidth}
\textbf{Step 1: Define the variable.} $T$ is the trial number of the first stated success, in trials, with support $\{1,2,\ldots\}$. $p=0.15$ and $1-p=0.85$.
\end{minipage}\par\medskip
\par\noindent\begin{minipage}{\linewidth}
\textbf{Step 2: Check assumptions.} The trials are binary, independent and have constant success probability; the process stops at the first success. Thus a geometric total-trial model applies.
\end{minipage}\par\medskip
\par\noindent\begin{minipage}{\linewidth}
\textbf{Step 3: Translate and calculate.} Use
\[P(T>m)=(1-p)^m.\]
Then
\[P(T>5)=(0.85)^5\approx 0.4437.\]
\end{minipage}\par\medskip
\par\noindent\begin{minipage}{\linewidth}
\textbf{Step 4: Check and interpret.} The result is in $[0,1]$; by-trial and after-trial probabilities sum to 1 because their events are complements. The endpoint is included for by and excluded for after. 
\end{minipage}\par\medskip
\secInicio[][sec:t2-c5-intermediate]
\Needspace{19\baselineskip}
\section{Advanced contextual practice}\label{sec:t2-c5-advanced}
\Needspace{14\baselineskip}
\subsection{Call budget and first sale}
Independent calls each sell a service with probability 0.25. Each call costs 3 USD. The business can fund at most 12 USD of calls.

\textbf{Questions/Tasks}
\begin{enumerate}
\item Calculate $P(T\leq4)$. Translate the 12 USD budget into a trial cutoff.
\end{enumerate}
\subsubsection*{Complete worked solution}
\par\noindent\begin{minipage}{\linewidth}
\textbf{Step 1: Define the variable.} $T$ is the trial number of the first stated success, in trials, with support $\{1,2,\ldots\}$. $p=0.25$ and $1-p=0.75$.
\end{minipage}\par\medskip
\par\noindent\begin{minipage}{\linewidth}
\textbf{Step 2: Check assumptions.} The trials are binary, independent and have constant success probability; the process stops at the first success. Thus a geometric total-trial model applies.
\end{minipage}\par\medskip
\par\noindent\begin{minipage}{\linewidth}
\textbf{Step 3: Translate and calculate.} Use
\[P(T\leq m)=1-(1-p)^m.\]
Then
\[P(T\leq4)=1-(0.75)^4\approx 0.6836.\]
\end{minipage}\par\medskip
\par\noindent\begin{minipage}{\linewidth}
\textbf{Step 4: Check and interpret.} The result is in $[0,1]$; by-trial and after-trial probabilities sum to 1 because their events are complements. The endpoint is included for by and excluded for after. $C=3T\leq12$ is equivalent to $T\leq4$. The budget is money; 4 is a call count. The probability is that a first sale arrives before the budget is exceeded.
\end{minipage}\par\medskip
\secInicio[][sec:t2-c5-advanced]
\Needspace{14\baselineskip}
\subsection{Failures before an insurance sale}
Independent insurance offers each sell with probability 0.1. Let $F$ count unsuccessful offers before the first sale.

\textbf{Questions/Tasks}
\begin{enumerate}
\item Calculate $P(T=4)$. Translate $F=3$ into total-trial notation, and contrast the event with exactly one sale among four offers.
\end{enumerate}
\subsubsection*{Complete worked solution}
\par\noindent\begin{minipage}{\linewidth}
\textbf{Step 1: Define the variable.} $T$ is the trial number of the first stated success, in trials, with support $\{1,2,\ldots\}$. $p=0.1$ and $1-p=0.9$.
\end{minipage}\par\medskip
\par\noindent\begin{minipage}{\linewidth}
\textbf{Step 2: Check assumptions.} The trials are binary, independent and have constant success probability; the process stops at the first success. Thus a geometric total-trial model applies.
\end{minipage}\par\medskip
\par\noindent\begin{minipage}{\linewidth}
\textbf{Step 3: Translate and calculate.} Use
\[P(T=t)=(1-p)^{t-1}p.\]
Then
\[P(T=4)=(0.9)^{3}(0.1)\approx 0.0729.\]
\end{minipage}\par\medskip
\par\noindent\begin{minipage}{\linewidth}
\textbf{Step 4: Check and interpret.} The result is in $[0,1]$; by-trial and after-trial probabilities sum to 1 because their events are complements. For the exact event, 3 failures precede success on trial 4. $F=3$ is $T=4$, sequence FFFS. Exactly one sale in four fixed offers also includes SFFF, FSFF and FFSF, so it is a binomial event with probability $\binom41(0.1)(0.9)^3=0.2916$, four times the waiting probability.
\end{minipage}\par\medskip
\secInicio[][sec:t2-c5-advanced]
\section{Closing reference: key ideas and common errors}
$T$ counts the successful trial as well as preceding failures. Its support starts at 1; $F=T-1$ starts at 0. Use $t-1$ failures in an exact probability and $m$ failures for after trial $m$. Distinguish first success from a fixed-trial success count, and state constant probability and independence.
\secInicio
\end{document}
